\chapter{Code Extraction and the Proof DSL}
\label{ch:extraction}

This chapter introduces two major extensions to the C$^4$ calculus:
\emph{code extraction} from verified proofs, and a \emph{pretty proof DSL}
that makes C$^4$ proofs readable as mathematical arguments.

\section{Code Extraction from Proofs}

\subsection{Motivation: Beyond F*}

F*~\cite{swamy2016} allows extracting OCaml, F\#, and C code from proofs
via its \texttt{effect} system.  C$^4$ extraction offers three material
advantages:

\begin{enumerate}
  \item \textbf{Same-language extraction.}  F* extracts to OCaml/F\#/C —
    languages unrelated to the verification target.  C$^4$ extracts to
    Python, the same language deppy verifies.  This eliminates the
    ``semantic gap'' between verified spec and extracted implementation.

  \item \textbf{Bidirectional extraction.}  F* provides only forward
    extraction (proof $\to$ code).  C$^4$ supports both:
    \begin{itemize}
      \item \emph{Forward}: proof term + OTerm spec $\to$ executable Python.
      \item \emph{Backward}: Python code + \texttt{@guarantee} annotation
            $\to$ C$^4$ proof obligation.
    \end{itemize}
    The backward direction is deppy's \texttt{@guarantee} mechanism,
    compiled to a C$^4$ $\Sigma$-type.

  \item \textbf{Direct denotational semantics.}  F* proofs are written
    against abstract specifications; the link to concrete code requires
    a separate extraction and erasure pass.  C$^4$ proofs are against
    OTerms — the actual denotational semantics of the source — so
    extraction is a direct projection, not a translation.
\end{enumerate}

\subsection{The Extraction Algorithm}

\begin{definition}[Code Extraction]
Given a proof $p : f \equiv g$ and an OTerm $f$, the extraction
$\mathsf{extract}(p, f)$ produces executable Python code $\mathsf{code}$
such that $\den{\mathsf{code}} = f$.
\end{definition}

The extraction proceeds by structural recursion on the proof term:

\begin{enumerate}
  \item \textbf{$\Sigma$-type projection.}  A proof of
    $\Sigma(x : \mathsf{Code}). \mathsf{spec}(x)$ extracts to $x$
    by first projection.

  \item \textbf{NatInduction.}  An induction proof
    \[
      \frac{\Gamma \vdash p_0 : f(0) \equiv g(0) \qquad
            \Gamma, n, \mathsf{ih} : f(n) \equiv g(n) \vdash p_s : f(n{+}1) \equiv g(n{+}1)}
           {\Gamma \vdash \mathsf{NatInd}(p_0, p_s) : \forall n.\; f(n) \equiv g(n)}
    \]
    extracts to either a recursive function (matching $f$) or an
    iterative loop (matching $g$), depending on which side is preferred.

  \item \textbf{Simulation.}  A bisimulation proof extracts the
    concrete implementation side, since the simulation relation
    guarantees behavioral equivalence.

  \item \textbf{CechGluing.}  A \v{C}ech gluing proof over fibers
    $\{U_i\}$ extracts a dispatch function:
    \begin{lstlisting}
def extracted(x):
    fiber = classify(x)
    if fiber == 'U0':
        return handle_U0(x)
    elif fiber == 'U1':
        return handle_U1(x)
    ...
    \end{lstlisting}
    Each handler is extracted from the corresponding local proof.

  \item \textbf{AbstractionRefinement.}  Both sides refine a common
    specification; extraction picks the preferred concrete implementation.

  \item \textbf{LoopInvariant.}  A loop invariant proof extracts a
    \texttt{while} loop with the invariant as an assertion:
    \begin{lstlisting}
def extracted(args):
    state = init(args)
    while not done(state):
        assert invariant(state)
        state = step(state)
    return output(state)
    \end{lstlisting}
\end{enumerate}

\subsection{OTerm to Python Compilation}

The core of extraction is the function $\mathsf{oterm\_to\_python}$
that compiles OTerms to Python expressions:

\begin{center}
\begin{tabular}{ll}
\toprule
OTerm constructor & Python output \\
\midrule
$\mathsf{OVar}(x)$ & \texttt{x} \\
$\mathsf{OLit}(v)$ & \texttt{repr(v)} \\
$\mathsf{OOp}(+, [a, b])$ & \texttt{(a + b)} \\
$\mathsf{OCase}(c, t, e)$ & \texttt{(t if c else e)} \\
$\mathsf{OLam}(x, b)$ & \texttt{lambda x: b} \\
$\mathsf{OFold}(\mathit{init}, \mathit{coll}, \mathit{step})$ &
  \texttt{functools.reduce(...)} or \texttt{for} loop \\
$\mathsf{OFix}(f, b)$ & recursive \texttt{def f(...): ...} \\
\bottomrule
\end{tabular}
\end{center}

\subsection{The \texttt{@verified} Decorator}

Extracted code is annotated with a \emph{verification certificate}:

\begin{lstlisting}
from cctt.proof_theory.extraction import verified

@verified(CORRECTNESS_PROOF, IMPL_OTERM, SPEC_OTERM)
def floyd_warshall(W, n):
    ...
\end{lstlisting}

The decorator attaches \texttt{\_c4\_certificate},
\texttt{\_c4\_verified}, and \texttt{\_c4\_proof\_strategy}
attributes to the function, enabling downstream consumers to
inspect the trust level without re-verifying.

\subsection{Bidirectional Extraction: The Key Advantage}

\begin{theorem}[Bidirectional Extraction Soundness]
Let $\mathsf{code}$ be Python source with annotation
$\texttt{@guarantee}(s)$.
\begin{enumerate}
  \item (Forward) If $p : f \equiv g$ is verified, then
    $\mathsf{extract}(p, f)$ produces code $c$ with
    $\den{c} = f$ and $c$ carries guarantee $s$.
  \item (Backward) $\mathsf{code\_to\_obligation}(\mathsf{code}, s)$
    produces a proof obligation $(f, \Sigma\text{-type})$ such that
    any valid proof $p : \den{\mathsf{code}} \in \Sigma\text{-type}$
    certifies that $\mathsf{code}$ satisfies $s$.
\end{enumerate}
\end{theorem}

This bidirectionality means the extraction and verification pipelines
form a closed loop: verified proofs produce guaranteed code, and
guaranteed code generates proof obligations that can be verified.


\section{The Pretty Proof DSL}

\subsection{Motivation}

Raw C$^4$ proof terms are compositional but verbose:

\begin{lstlisting}
NatInduction(
    base_case=Assume('base_eval', base_lhs, base_rhs),
    inductive_step=Trans(
        Cong('*', [Refl(var('n')), Assume('ih', ...)]),
        AxiomApp('fold_spec', 'backward', FACT_ITER),
    ),
    variable='n',
)
\end{lstlisting}

The DSL module provides two interfaces for writing cleaner proofs:

\subsection{The ProofBuilder}

\texttt{ProofBuilder} is a fluent API for proof construction:

\begin{lstlisting}
from cctt.proof_theory.dsl import ProofBuilder

pb = ProofBuilder(lhs=FACT_REC, rhs=FACT_ITER)
base = pb.refl(lit(1))
ih = pb.assume('ih', FACT_REC, FACT_ITER)
step = pb.calc_trans(
    pb.congr('*', pb.refl(var('n')), ih),
    pb.axiom('fold_spec', 'backward', FACT_ITER),
)
pb.by_induction(on='n', base=base, step=step)
result = pb.verify()   # returns VerificationResult
\end{lstlisting}

\subsection{The \texttt{@c4\_proof} Decorator}

\begin{lstlisting}
@c4_proof(lhs=FACT_REC, rhs=FACT_ITER)
def factorial_equiv(pb):
    """fact_rec == fact_iter by induction on n"""
    base = pb.refl(lit(1))
    ih = pb.assume('ih', FACT_REC, FACT_ITER)
    step = pb.calc_trans(
        pb.congr('*', pb.refl(var('n')), ih),
        pb.axiom('fold_spec', 'backward', FACT_ITER),
    )
    pb.by_induction(on='n', base=base, step=step)

# factorial_equiv is a ProofBuilder; call .build() or .verify()
\end{lstlisting}

\subsection{Calculational Proofs}

The DSL supports \texttt{calc}-style chains (analogous to Lean's \texttt{calc}):

\begin{lstlisting}
pb.calc_chain(
    (term_b, proof_a_eq_b),    # a == b by proof_a_eq_b
    (term_c, proof_b_eq_c),    # b == c by proof_b_eq_c
    (term_d, proof_c_eq_d),    # c == d by proof_c_eq_d
)
\end{lstlisting}

This compiles to a \texttt{RewriteChain} proof term.

\subsection{The ProofScript}

For longer proofs, \texttt{ProofScript} provides an imperative interface:

\begin{lstlisting}
script = ProofScript(lhs=FW_OTERM, rhs=SPEC_OTERM)
script.induction('k')
script.base_case(fw_base_proof)
script.inductive_step(
    script.calc(
        (D_k_plus_1, unfold_proof),
        (relaxed, bellman_proof),
        (spec_k_plus_1, induction_hypothesis),
    )
)
proof = script.qed()
\end{lstlisting}


\section{Python-Specific Proof Patterns}

The \texttt{python\_patterns} module provides 10 pre-built proof templates
for common Python equivalences:

\begin{center}
\begin{tabular}{lll}
\toprule
Pattern & Equivalence & Strategy \\
\midrule
\texttt{list\_comp} & Comprehension $\equiv$ loop-append & ListInduction \\
\texttt{generator} & Generator $\equiv$ materialized list & Simulation \\
\texttt{mutation\_isolation} & deepcopy safety & Simulation \\
\texttt{exception} & try/except $\equiv$ if/else & CasesSplit \\
\texttt{decorator} & Transparent decorator & Extensionality \\
\texttt{context\_manager} & with $\equiv$ try/finally & Simulation \\
\texttt{map\_filter} & map/filter compose & ListInduction \\
\texttt{reduce\_loop} & reduce $\equiv$ for-loop & Reflexivity \\
\texttt{dict\_comp} & Dict comprehension $\equiv$ loop & ListInduction \\
\texttt{sorted\_key} & sorted(key=f) $\equiv$ manual sort & Abstraction \\
\bottomrule
\end{tabular}
\end{center}

Usage:
\begin{lstlisting}
from cctt.proof_theory.python_patterns import by_pattern
inst = by_pattern('list_comp', collection=OVar('xs'),
                  transform=OLam('x', OOp('+', (OVar('x'), OLit(1)))))
result = inst.verify()
\end{lstlisting}


\section{Spec Compilation: @guarantee to $\Sigma$-types}

\subsection{The Compilation Pipeline}

A deppy \texttt{@guarantee} annotation like:

\begin{lstlisting}
@guarantee("returns a sorted list with no duplicates")
def unique_sorted(items: list[int]) -> list[int]:
    ...
\end{lstlisting}

compiles to the C$^4$ $\Sigma$-type:

\[
  \Sigma(\mathit{result} : \mathsf{List}).\ 
    \underbrace{(\forall i.\ \mathit{result}[i] \le \mathit{result}[i{+}1])}_{\text{structural (Z3)}}
    \wedge
    \underbrace{(\forall i \ne j.\ \mathit{result}[i] \ne \mathit{result}[j])}_{\text{structural (Z3)}}
\]

The compiler:
\begin{enumerate}
  \item Parses the NL spec into keyword-triggered predicates.
  \item Classifies predicates as \emph{structural} (Z3-decidable) or
        \emph{semantic} (oracle-judgeable).
  \item Generates proof obligations for each predicate.
  \item Produces a proof skeleton with \texttt{Z3Discharge} nodes for
        structural obligations and \texttt{Assume} nodes for semantic ones.
\end{enumerate}

\subsection{Structural vs.\ Semantic Predicates}

\begin{definition}[Hybrid Spec Decomposition]
A guarantee $s$ decomposes as $s = s_{\mathrm{str}} \wedge s_{\mathrm{sem}}$ where:
\begin{itemize}
  \item $s_{\mathrm{str}}$ is the conjunction of Z3-decidable predicates
    (sortedness, uniqueness, bounds, etc.).
  \item $s_{\mathrm{sem}}$ is the conjunction of predicates requiring
    oracle judgment (correctness, optimality, safety, etc.).
\end{itemize}
\end{definition}

This decomposition mirrors the fiber structure of C$^4$: structural predicates
live in the decidable fiber, semantic predicates in the oracle fiber.


\section{Example: Floyd--Warshall Extraction}

The Floyd--Warshall proof from Chapter~\ref{ch:proof-examples} (if available)
demonstrates end-to-end extraction:

\begin{lstlisting}
from cctt.proof_theory.extraction import extract_from_proof

extracted = extract_from_proof(
    proof=CORRECTNESS_PROOF,
    lhs=FW_OTERM,
    rhs=SPEC_OTERM,
    function_name='floyd_warshall',
    params=['W', 'n'],
)

print(extracted.with_guarantee())
# @guarantee("equivalent to shortest_path_spec by induction")
# def floyd_warshall(W, n):
#     ...

print(extracted.certificate.trust_level())
# 'CONDITIONAL' (uses IH assumption)
\end{lstlisting}

The extracted code carries its proof certificate, enabling
downstream consumers to inspect the verification status
without re-running the proof checker.
